\subsection{Unit-root eighth moments for the terminal \((1,2)\)}
\label{proof:Q5-C050}
Use the original finite probability measure, \(A=J+K\), \(T=J-K\),
\(g=K\one\), \(p=1+g\), \(q=1-g\),
\(r=Kg\), \(s=Kr\), \(t=Ks\), and
\(d=\|g\|^2\), \(\delta_j=\|K^j\one\|^2\).
Write \(\theta=13/32\), \(\kappa=19/32\),
\(\tau=d-\delta_2\), and \(\gamma=d-\delta_2/2\).
Lemma~\ref{q3:lem:spectral}, skew-adjointness, and the row bound give
\[
 0\le d\le1/3,\qquad d^2\le\delta_2\le\theta d,
 \qquad\delta_{j+1}\le\theta\delta_j,\qquad |Kv|\le\|v\|.
\]
All connector edges point from the terminal toward the receiving root,
whose original unit arm points outward. The terminal unit arm points
inward and its two-edge arm outward. Thus the actual message is
\begin{equation}\tag{Q50.0.1}\label{eq:Q50:0.1}
 F=qf_2,\qquad X_L=T^LF,\qquad H_L=pX_L.
\end{equation}
We shall prove the stronger estimate \(\E H_L^8\le1-d\) for \(L\ge2\);
in particular,
\begin{equation}\tag{Q50.0.2}\label{eq:Q50:0.2}
 \E H_L^8\le1-d/50\qquad(L\ge2).
\end{equation}
The case \(L=1\) is Section~\ref{proof:Q5-C049}.
When \(d=0\), leaf integration makes every tree message one.
Assume \(d>0\) below.

\paragraph{A common unit-root payment.}
Suppose a nonnegative message satisfies
\begin{equation}\label{eq:unit-ordinary:positive-square}
 H\le B\left(1+\frac{Z-a}{2B}\right)^2,
 \qquad |Z|\le U,\quad \E Z\le-\alpha,\quad \E Z^2\le v.
\end{equation}
Here \(B>0\), \(a,\alpha\ge0\), \(\alpha^2\le v\le U^2\),
and \(2B-a>U>0\).
Apply Lemma~\ref{lem:unit-square-hermite} with exponent eight,
the variable \(Z/\sqrt B\), and center \((B-a/2)/\sqrt B\).
The compatible double contact and upper probability are
\[
 \ell=-\frac{v+\alpha U}{U+\alpha},\qquad
 \pi=\frac{v-\alpha^2}{U^2+2\alpha U+v}.
\]
Consequently,
\begin{equation}\label{eq:unit-ordinary:compatible-sixteen}
 \E H^8\le
 \frac{(1-\pi)(B-a/2+\ell/2)^{16}
             +\pi(B-a/2+U/2)^{16}}{B^8}.
\end{equation}
At \(v=U^2\), the same Hermite quadratic has contacts \(-U,U\),
with upper probability \((U-\alpha)/(2U)\), so the displayed formula
also holds.
This uses the moments of the same $Z$.

Here are five continuous payments for this formula. Put
\(c=51/64\), \(\beta=89/64\), \(k=19/36\), and
\[
 v_3(d)=\theta(1-cd)^2+\theta^2\kappa^2
                  -2\kappa^2\theta d+\kappa^2d^2.
\]
The table writes \(U=t u\), \(\alpha=d a_0\), \(v=d v_0\),
where \(d=t^2\). Every row implies \(\E H^8\le1-d\).
\begin{table}[htbp]\centering\small
\caption{Joint moment data for the unit \((1,2)\) terminal.}
\label{tab:unit-ordinary:two-data}
\setlength{\tabcolsep}{3pt}
\begin{tabular}{@{}cllllll@{}}\toprule
Row & \(B\) & \(a\) & \(a_0\) & \(u\) & \(v_0\) & Actual \(d\)\\\midrule
I & \(1-cd\) & \(0\) & \(\kappa\) & \(1-cd\) & \(4/27\) & \([0,1/4]\)\\
II & \(1-\beta d\) & \(0\) & \(1-cd\) & \(\frac{43}{40}(1-\frac45d)\) & \(v_3(d)\) & \([0,1/3]\)\\
III & \(1-\beta d\) & \(d(1-d)\) & \(\frac{777}{1024}-d\) & \(1\) & \(\theta(1-\frac{153}{128}d)\) & \([0,1/3]\)\\
IV & \(1-3kd\) & \(kd(1-d)\) & \(k(1-2d)\) & \(1\) & \(1/4\) & \([0,1/3]\)\\
V & \(1-3\kappa d\) & \(\kappa d(1-d)\) & \(\kappa(1-2d)\) & \(1\) & \(1/3\) & \([0,1/3]\)\\\bottomrule
\end{tabular}
\end{table}

The single polynomial certificate for each row can be
reconstructed as follows:
\[
 \begin{gathered}
 C=B-a/2,\quad D=u+ta_0,\quad E=u^2+2tua_0+v_0,\quad
 N=v_0+ta_0u,\quad R=v_0-t^2a_0^2,\\
 \mathcal G(t)=\frac{
  2^{16}D^{14}EB^8(1-t^2)
 -(2CD-tN)^{16}-RD^{14}(2C+tu)^{16}}{t^2}.
 \end{gathered}
\]
The numerator is divisible by \(t^2\), and \(\mathcal G(0)>1\) in
every row. The denominator of the bound
\eqref{eq:unit-ordinary:compatible-sixteen} is
\(2^{16}D^{14}EB^8>0\).
On the certificate intervals in the next table, every Bernstein
coefficient of \(\mathcal G/\mathcal G(0)\) exceeds \(1/4000\).
The six guard polynomials
\[
 B,\quad u,\quad a_0,\quad v_0-t^2a_0^2,
 \quad u^2-v_0,\quad 2B-a-tu
\]
each have Bernstein coefficients greater than \(1/32\) on the same
interval. In particular all the Hermite hypotheses hold strictly when
\(t>0\). The coefficients follow from
\eqref{eq:bernstein-coefficients}.
\begin{table}[htbp]\centering\small
\caption{Whole-interval certificates for Table~\ref{tab:unit-ordinary:two-data}.}
\label{tab:unit-ordinary:two-certificates}
\begin{tabular}{@{}cccc@{}}\toprule
Row & Interval for \(t\) & Degree of \(\mathcal G\) & Coefficient count\\\midrule
I & \([0,1/2]\) & \(76\) & \(77\)\\
II & \([0,7/12]\) & \(93\) & \(94\)\\
III--V & \([0,7/12]\) & \(107\) & \(108\) each\\\bottomrule
\end{tabular}
\end{table}
Since \((7/12)^2>1/3\), these rational intervals cover every actual
\(d\)-domain. Lemma~\ref{lem:bernstein-certificate} now proves the five
payments, including the limiting endpoint \(d=0\).

\paragraph{The finite terminal envelope.}
The arm identity and positivity give
\[
 qf_2=(1+r/2)^2-(g+r/2)^2,\qquad
 TF\le q-\gamma-s.
\]
Put \(B_0=-K\), \(t_j=B_0^j\one\),
\(S'_L=\sum_{j=0}^Lt_j\), and \(\bar f_n=T^n\one\).
Let \(\beta_1=\beta_2=0\) and
\(\beta_j=\sum_{i=1}^{\lfloor(j-1)/2\rfloor}(-1)^{i-1}\delta_i\)
for \(j\ge3\); alternating sums give \(\tau\le\beta_j\le d\).
Apply \(T^{L-1}\) to the one-step bound above. The finite convolution
\eqref{eq:Q188:conv}, or direct finite Duhamel expansion, gives for
every \(L\ge1\)
\begin{equation}\label{eq:unit-ordinary:two-finite}
 \begin{gathered}
 X_L\le\mathcal W_L:=S'_L+t_{L+2}-\gamma\bar f_{L-1}
                    -\sum_{j=3}^L\omega_j\bar f_{L-j},\\
 \omega_j=\begin{cases}\beta_j,&j\text{ even},\\
                         \beta_{j+2},&j\text{ odd},\end{cases}
 \qquad\omega_j\ge\tau.
 \end{gathered}
\end{equation}
Only nonnegative indices occur. Indeed
\(T^Lr=t_{L+2}-d\bar f_{L-1}+\delta_2\bar f_{L-3}
-\delta_3\bar f_{L-5}+\cdots\), a finite identity.
All retained path losses are nonnegative.

\paragraph{Two and three connector edges.}
For \(L=2\), \(X_2\le(1-\gamma)q+r+t\).
At \(d\le1/4\), weaken \(\gamma\ge cd\) and use
\(Z=r+t\). Its mean is \(-d+\delta_2\le-\kappa d\), and
\[
 \E Z^2=\delta_2-2\delta_3+\delta_4\le4d/27,
 \qquad |Z|^2\le d-2\delta_2+\delta_3
                  \le d(1-2cd).
\]
The first bound follows from
\(4/27-x(1-x)^2=(x-1/3)^2(4/3-x)\ge0\).
For the row bound use \(Z=K(g+s)\), center the row, and note that
\(g+s\) is centered. Since
\(\sqrt{1-2cd}\le1-cd\), these are row I.
For \(d\ge1/4\), \(TF\le2\E F=2-4d\le1\), hence
\(H_2\le(2-4d)pq\le1\) and
\(\E H_2^8\le\E(pq)^8\le1-d\).

For \(L=3\), expand \eqref{eq:unit-ordinary:two-finite} to obtain
\[
 X_3\le(1-\beta d)q+Z_3,\qquad
 Z_3=(1-cd)r-s-K^4g-\kappa dg,
 \qquad\E Z_3=-d(1-cd).
\]
Writing \(A_0=1-cd\), the row and spectral bounds are
\[
 |Z_3|^2\le dD_3(d),\quad
 D_3=A_0^2+4/27-\kappa^2d,
 \qquad \E Z_3^2\le d v_3(d).
\]
For the row estimate use
\(Z_3=(-K)(-A_0g+r+t+\kappa d)\): its squared predecessor norm
is at most \(A_0^2d+4d/27-\kappa^2d^2\), since
\(\E(r+t)=-\tau\) and \(\tau\ge\kappa d\).
For its second moment the exact spectral integrand is
\(A_0^2x+[x(1-x)-\kappa d]^2\).
Its derivative is at least \(A_0^2-2\kappa d>0\), so its value
at \(x=\theta\) gives \(v_3\).
Finally \(D_3\le(43/40)^2(1-4d/5)^2\): the three coefficients
of the difference are
\(1849/1600-31/27\),
\(1993/1024-1849/1000\), and
\(1849/2500-2601/4096\), all positive. Thus row II applies.

\paragraph{Four connector edges.}
Set \(R_L=\sum_{j=2}^Lt_j+t_{L+2}\).
Expansion of the same finite envelope gives
\[
 X_4\le(1-\tau-\gamma)q+Z_4-d(1-\gamma),
 \qquad Z_4=R_4-\gamma r+\gamma s.
\]
Here \(Z_4=B_0Q(B_0)g'\), \(g'=B_0\one\),
\(Q=P_3-\gamma(1+z)\), \(P_3=1+z+z^2+z^4\).
At \(z=i\sqrt x\),
\[
 |P_3|^2=1-x+x^2(3-2x+x^2)\le1,
 \quad
 |Q|^2\le1-2\gamma+(1+\theta)\gamma^2\le1.
\]
The first inequality uses
\(1-3x+2x^2-x^3\ge1451/32768>0\).
Consequently \(|Z_4|\le\sqrt d\), and
\[
 \E Z_4^2\le\theta d(1-153d/128),\qquad
 \E Z_4\le-(777/1024)d+d^2.
\]
For the variance use \(\gamma\ge cd\),
\((1+\theta)cd\le1/2\); for the mean use
\(\tau+\delta_3\ge d[1-u+u^2]\ge(777/1024)d\),
where \(u=\delta_2/d\le\theta\) and
\(\delta_3\ge\delta_2^2/d\).
Weakening only positive path losses gives
\(X_4\le(1-\beta d)q+Z_4-d(1-d)\), namely row III.

\paragraph{Five edges and all longer connectors.}
At \(L=5\), retain the four losses already present in
\eqref{eq:unit-ordinary:two-finite}. Expanding
\(\bar f_4=(1-d)q-d+r-s+t\) gives
\[
 X_5\le[1-(3-d)\tau]q+Z_5-\tau(1-d),
 \qquad Z_5=R_5-2\tau r+\tau s-\tau t.
\]
Its spectral polynomial is
\(Q=P_4-\tau(2+z+z^2)\), \(P_4=1+z+z^2+z^3+z^5\).
The identities
\[
 |P_4|^2\le1-x,\quad
 \Re[P_4\overline{(2+z+z^2)}]=2-2x+x^3\ge19/16,
 \quad |2+z+z^2|^2\le4
\]
and \(\tau\le2/9\) imply \(|Q|^2\le|P_4|^2\).
Thus \(|Z_5|\le\sqrt d\), \(\E Z_5^2\le d/4\), and
\(\E Z_5=-\tau+2\tau d-\tau\delta_2
\le-\tau(1-2d)\).
Since \((3-d)\tau\ge3kd\), row IV applies.

For \(L\ge6\), retain the last four losses, whose indices are all
at least three. They give
\begin{equation}\label{eq:unit-ordinary:two-long}
 X_L\le(1-3\tau)q+Z_L-\tau(1-d),
 \qquad Z_L=R_L-2\tau r+\tau s.
\end{equation}
For \(m=L-1\ge5\), put
\(P_m=G_m+z^{m+1}\), \(G_m=\sum_{j=0}^{m-1}z^j\), and
\(Q_m=P_m-\tau(2+z)\). Direct multiplication gives
\[
 |P_m|^2=
 \frac{1+x^m(1-x+x^2)-2\Re[z^m(1-x-z)]}{1+x}.
\]
For \(m\equiv0,1\pmod4\), the subtracted terms are respectively
\(2x^{m/2}(1-x)\) and \(2x^{(m+1)/2}\); both dominate
\(x^m(1-x+x^2)\). For the other residues the numerator is at most
\(1+2x^3+x^6\) or \(1+2x^4+x^7\).
These estimates yield \(|P_m|\le1\), \(x|P_m|^2\le1/3\).
For the latter bound in the last two cases it suffices to use
\(2x+6x^4+3x^7\le2269/2304<1\).
Also
\(\Re[P_m\overline{(2+z)}]\ge2-3x/(1-x^2)\ge154/285\).
Together with \(\tau(4+x)\le47/48<308/285\), this proves
\(|Q_m|^2\le|P_m|^2\).
Therefore
\[
 |Z_L|\le\sqrt d,\qquad \E Z_L^2\le d/3,
 \qquad\E Z_L\le-\tau(1-2d).
\]
For the last bound \(\E R_L=-\beta_{L+3}\) when \(L\) is even,
and \(-\beta_{L+1}\) when \(L\) is odd.
Replacing \(\tau\) by \(\kappa d\) only on the positive path
factors in \eqref{eq:unit-ordinary:two-long} gives row V.
In each of the five cases, multiplying the nonnegative envelope
\(X_L\le Bq+Z-a\) by \(p\) and applying AM--GM to
\(Bp\) and \(Bq+Z-a\) proves
\eqref{eq:unit-ordinary:positive-square}. This completes the proof.

\paragraph{A scalar estimate used by the neighbouring terminal family.}
We retain the elementary bound
\begin{equation}\tag{Q50.2.1}\label{eq:Q50:2.1}
 (1+x)^{16}\le1+16x+256x^2\qquad(|x|\le1/7).
\end{equation}
On the negative half-interval Taylor's quadratic coefficient is at most
\(120\); on the positive half the quotient is increasing, and at
\(1/7\) the assertion is \(8^{16}\le417\,7^{14}\), whose margin
is \(1344044667377>0\).

Finally, actual equal-pair factors \(E_r,E_t\le1\) satisfy
\(0\le E_rpT^L(E_tqf_2)\le H_L\) by positivity. Full reversal
replaces \(K\) by \(-K\) and reverses every arm and connector.
Together with Section~\ref{proof:Q5-C049}, the catalogue for all
positive lengths retains
\begin{equation}\tag{Q50.8.1}\label{eq:Q50:8.1}
 \E H_L^8\le1-d/500\qquad(L\ge1).
\end{equation}
